\section{Preliminaries}
\label{sec:prelim}

For an integer $t \ge 1$ we denote $[t]:=\{1,2,\ldots,t\}$. The statistical (or total variation) distance between two distributions $\mu,\mu'$ is
given by $\sd(\mu,\mu')=\sum_{x} |\mu(x)-\mu'(x)|$. The expression $\delta_{i,j}$ is equal to $1$ if $i=j$ and is equal to $0$ otherwise.



\paragraph{Group action and uniformity}
A group $G$ \emph{acts} on a set $X$ if there is a homomorphism from $G$ to the
permutation group on $X$. That is, each $g \in G$ is a permutation on $X$,
and $(gh)(x)=g(h(x))$ for all $g,h \in G, x \in X$. We denote by $U_G$
the uniform distribution over $G$. For a distribution $\mu$ over $G$ we denote
by $g \sim \mu$ a random element chosen according to $\mu$. We abbreviate $\Pr_{g \in G}[\cdot] = \Pr_{g \sim U_G}[\cdot]$.

We recall some definitions from the
introduction: a distribution $\mu$ over $G$ is \emph{$X$-uniform} if
$$
\Pr_{g \sim \mu}[g(x)=y] = \Pr_{g \in G}[g(x)=y]
$$
for all $x,y \in X$; and a distribution $\mu$ is
\emph{almost $X$-uniform with error $\eps$} if
$$
\left| \Pr_{g \sim \mu}[g(x)=y] -
\Pr_{g \in G}[g(x)=y] \right| \le \eps
$$
for all $x,y \in X$. A family $S \subset G$ is \emph{$X$-uniform} (\emph{almost
$X$-uniform}, accordingly) if the uniform distribution over $S$ is such. We will need
some basic facts in linear algebra, geometry and representation theory,
which are presented below.

\paragraph{Linear algebra}
Let $A=(a_{i,j})$ be a complex matrix. The $L_{\infty}$ norm of $A$ is $\|A\|_{\infty}=\max
|a_{i,j}|$ (not to be confused with the operator norm of $A$ on $L_{\infty}$, which is not used in this paper).
The Frobenius norm of $A$ is $\|A\|_{\Fr} = \sqrt{\sum
|a_{i,j}|^2}$. For every $A$, $\|A\|_{\infty} \le \|A\|_{\Fr}$. A matrix
$A$ is unitary if $A \overline{A^t} = I$. The Frobenius
norm of a matrix is invariant under any unitary
basis change. That is, if $U$
is unitary then $\|U^{-1} A U\|_{\Fr} = \|A\|_{\Fr}$. The tensor product
of a $d_1 \times d_1$ matrix $A_1$ and a $d_2 \times d_2$ matrix $A_2$,
denoted $A_1 \otimes A_2$, is a $(d_1 d_2) \times (d_1 d_2)$ matrix, whose
entries are given by $(A_1 \otimes A_2)_{(i,i'),(j,j')} = (A_1)_{i,j}
(A_2)_{i',j'}$. The tensor product of unitary matrices is
also unitary.

\paragraph{Geometry}
Let $X=\{x_1,\ldots,x_N\}$ be a set of points in $\R^d$. The \emph{convex hull} of $X$ is the set of points contained in the minimal convex
set containing $X$; equivalently, it is the set of all points 
\[
\Bigl\{\sum
\lambda_i x_i: \lambda_1,\ldots,\lambda_N \ge 0, \;\sum \lambda_i=1\Bigr\}\,.\]

\begin{fact}[Carath\'{e}odory theorem~\cite{Caratheodory}]\label{fact:caratheodory}
Let $X$ be a finite set of points in $\R^d$, and let $y$ be a point
in the convex hull of $X$. Then there exists a subset $Y \subset X$
of size $|Y| \le d+1$ such that $y$ is in the convex hull of $Y$.
\end{fact}

Any hyperplane $H$ partitions a set $X$ of points into two sets: if
$H=\{x:\ip{a,x}=b\}$ then the sets are $\{x \in X:\ip{a,x} \ge b\}$
and $\{x \in X:\ip{a,x}<b\}$. We need the following bound on the maximal
number of ways
a set can be partitioned by hyperplanes.\footnote{A quick way to prove a slightly weaker estimate
is as follows: the $VC$-dimension~\cite{VC71} of
halfspaces is $d+1$. Hence by the VC-dimension
theorem~\cite{VC71,Sauer72,Shelah72} the number of partitions
is at most $\sum_{i=0}^{d+1} {N \choose i} \le N^{d+1}$.}
\begin{fact}[Harding~\cite{Harding67}]
\label{fact:partition_hyperplane}
Let $X$ be a set of $N$ points in $\R^d$. The number
of different ways to partition $X$ into two sets by a hyperplane is
at most $\sum_{i=0}^d {{N-1} \choose i}  \leq N^{d}$.
\end{fact}

\paragraph{Representation theory}
Let $G$ be a finite group. A representation of $G$ (over $\C$) is a
homomorphism $R:G \to \GL(d,\C)$. That is, $R(g)$ for $g \in G$ is a
$d \times d$ nonsingular
complex matrix, and $R(gh) = R(g) R(h)$ for every $g,h \in
G$. The dimension of the representation $R$ is $d$. Two representations
$R,R'$ of $G$ of dimension $d$ are \emph{equivalent} if there exists an
invertible matrix $A$ such that $R'(g) = A^{-1} R(g) A$ for all $g \in
G$. This is denoted as $R \equiv R'$.

A representation $R$ is \emph{unitary}
if $R(g)$ is unitary for all $g \in G$.
\begin{fact}
Any representation of $G$ is equivalent to a unitary representation.
\end{fact}
We will restrict our attention only to unitary representations in this
paper. We note that if $R,R'$ are unitary and equivalent, then there
exists a unitary matrix $A$ such that $R'(g) = A^{-1} R(g) A$.

Let $G$ be a group which acts on a set $X$, that is, $g:X \to X$ is
a permutation of $X$ for all $g \in G$, and $(gh)(x)=g(h(x))$ for all
$g,h \in G$ and $x \in X$. The associated representation $R_X$ maps each
$g \in G$ to the permutation matrix it induces on the set $X$. That is,
$R_X(g)$ is an $|X| \times |X|$ matrix, indexed by $x,y \in X$, defined
as $(R_X(g))_{x,y}=1$ if $g(x)=y$ and $(R_X(g))_{x,y}=0$ otherwise. Note
that $R_X$ is always a unitary representation.

The sum of two representations $R_1:G \to \GL(d_1,\C)$ and $R_2:G \to
\GL(d_2,\C)$ is a representation $R:G \to \GL(d_1+d_2,\C)$ where $R(g)$ is
defined as a block diagonal matrix with two blocks, given by $R_1(g)$ and
$R_2(g)$. For $e \ge 1$ let $e R := R+\cdots+R$ where the sum is over $e$
copies of $R$. A representation $R$ is \emph{reducible} if it is equivalent
to the sum of two representations. Otherwise, the representation
$R$ is \emph{irreducible}. We summarize a few basic properties of
representations below. For details we refer the reader to any standard
book on representation theory, \eg,~\cite{RepresentationTheory}.

\begin{fact}[Maschke's theorem]\label{fact:maschke}
Any representation $R$ of $G$ is equivalent to a sum of irreducible
representations $R \equiv e_1 R_1+\cdots+e_ tR_t$, where $R_1,\ldots,R_t$
are nonequivalent irreducible representations, and $e_i \ge 1$ is the
multiplicity of $R_i$. We have $\sum e_i \dim(R_i) = \dim(R)$.
\end{fact}

\begin{fact}[Schur's lemma]\label{fact:schur}
Let $R$ be a unitary irreducible representation
of $G$ of dimension $d$. Then for any $i,j,k,\ell \in [d]$,
$$
\frac{1}{|G|} \sum_{g \in G} R(g)_{i,j}
\overline{R(g)_{k,\ell}} = \frac{1}{d} \delta_{i,k} \delta_{j,\ell}\,.
$$
Let $R',R''$ be two non-equivalent unitary irreducible
representations of $G$ of dimensions $d',d''$.
Then for any $i,j \in [d']$ and $k,\ell \in [d'']$,
$$
\frac{1}{|G|} \sum_{g \in G} R'(g)_{i,j} \overline{R''(g)_{k,\ell}} = 0\,.
$$
\end{fact}

The trivial representation is given by $\one(g)=1$ for all $g \in G$. An
immediate corollary of Schur's lemma is that
for every representation $R$ which is irreducible
and nontrivial, we have $\sum_{g \in G} R(g)=0$.

The group algebra $\C[G]$ is the linear space of functions $\mu:G \to
\C$. It is often written as $\mu = \sum_{g \in G} \mu(g) \cdot g$. Note
that the distributions over $G$ form
a subset of $\C[G]$. For $\mu \in \C[G]$
and a representation $R$ of $G$, let
$R(\mu) := \sum_{g \in G} \mu(g) R(g)$. If $\mu$ is a distribution,
this is equivalent to $R(\mu) = \E_{g \sim \mu}[R(g)]$. Note that in this
notation, the statistical distance between two distributions $\mu,\mu'$ on
$G$ is $\sd(\mu,\mu')=\|\mu-\mu'\|_1$.


